\section{Positive channels at equality}\label{sec:positive}

We prove that the limiting posterior is uniform when
\(0<\lambda<1\) and \(d\lambda^2=1\). Throughout this section,
\(m\) has the limiting fixed-point law supplied by Proposition~\ref{fnd:limit},
\(v=m-1\), and
\[
 \mu=\E x,\qquad \nu=\E x^2,\qquad \eta=\E y,\qquad
 \alpha=\frac{\E[D(D-1)]}{d^2}.
\]
Thus \(\alpha=1-1/d\) for the regular model and \(\alpha=1\) for the
Poisson model. A subscript \(e\) means evaluation at the edge message
\(m_e=1+\lambda v\). We argue by contradiction, assuming \(\mu>0\).

\subsection{An entropy identity with an additive correction}

For a positive message of coordinate average one, define
\begin{equation}\label{pos:functionals}
 \begin{aligned}
 I&=\av{m\log m},& J&=\frac12\av{v\log m},\\
 u_i&=\log m_i-\av{\log m},& B&=\frac12\av{m u^2}.
 \end{aligned}
\end{equation}
Here and below logarithms are natural. The expected additivity of $J$
below is the symmetric-channel form of the symmetrized-entropy recursion
of Formentin and K\"ulske~\cite[Proposition~3.3, equation~(15)]{FormentinKulske2009Entropy}:
their functional equals $2J$ in our normalization. We prove the identities
needed here, including the centered-log correction $B$.
Lemma~\ref{fnd:log-integrability} gives strictly positive fixed-point
coordinates and integrability of these functions: their coordinate factors
are bounded, and their logarithmic factors have degree at most two.

Under a fixed input spin \(i\), tilting and permutation symmetry give
\[
 \E[u_{e,i}\mid S=i]=\E\av{m_eu_e}=2\E J_e,
 \qquad
 \E[u_{e,i}^2\mid S=i]=\E\av{m_eu_e^2}=2\E B_e.
\]
Centered log coordinates add under products: the common normalizing
constant disappears on centering. Conditional independence of the branches
therefore gives
\begin{equation}\label{pos:additive}
 \E J=d\E J_e,\qquad
 \E B=d\E B_e+2\alpha(\E J)^2.
\end{equation}
For clarity, the off-diagonal term in the second formula, before division
by two, is \(4\E[D(D-1)](\E J_e)^2\).

For a product of \(n\) edge experiments write \(Z_n\) for its likelihood
normalizer relative to their independent marginal laws, and set
\[
 H_n=\E_{\mathrm{prod}}[Z_n\log Z_n],\qquad H=\E H_{D}.
\]
This is a relative entropy and is nonnegative. Expanding the logarithm of
the normalized output likelihood gives
\[
 \E I=d\E I_e-H.
\]
Each branch retains its original marginal distribution after the joint
likelihood tilt, which justifies the individual terms in this expansion.
The combination
\begin{equation}\label{pos:functional-F}
 F=3I-2J+\frac45B
\end{equation}
therefore satisfies the exact identity
\begin{equation}\label{pos:entropy}
 \E(F-dF_e)=-3H+\frac85\alpha(\E J)^2.
\end{equation}

The useful feature of this combination is its positive radial defect:
for every positive message of coordinate average one,
\begin{equation}\label{pos:radial}
 F(m)-\lambda^{-2}F(1+\lambda v)
 \ge\frac{1-\lambda^2}{2}x^2\qquad(0<\lambda<1).
\end{equation}
Lemma~\ref{rad:lemma} proves this pointwise inequality by radial
differentiation and explicit polynomial sign decompositions.
Since $d=\lambda^{-2}$, its expectation and
Equation~\eqref{pos:entropy} imply
\begin{equation}\label{pos:comparison}
 \frac{1-\lambda^2}{2}\nu+3H
 \le\frac85\alpha(\E J)^2.
\end{equation}
We will contradict this inequality by bounding $\E J$ above and $H$
below. The additive identity for $J$ also restricts the possible values
of $\mu,\nu,\eta$; we establish those restrictions first.

\subsection{Moment information from the additive identity}

The cubic identity in Equation~\eqref{fnd:cubic}, together with
$\av{v}=0$ and $\av{v^2}=2x$, gives the additional moment relation
\begin{equation}\label{pos:coordinate-identities}
 \av{v^4}=3x\av{v^2}+2y\av{v}=6x^2.
\end{equation}
We also use $|y|\le x^{3/2}$ from
Equation~\eqref{fnd:invariant-bounds}.

\begin{lemma}\label{pos:moment-bounds}
Put
\[
 c_0=\frac{37}{27},\qquad q=\frac{c_0}{2+\lambda},\qquad
 c=c_0+2\lambda q.
\]
Then
\begin{equation}\label{pos:moment-estimates}
 0\le\E J\le\mu-\lambda q\nu,
 \qquad c\nu\le\eta\le\sqrt{\mu\nu}.
\end{equation}
If \(s=1\) in the regular model and \(s\) is the survival probability in
the Poisson model, then
\begin{equation}\label{pos:moment-region}
 \frac{\mu^2}{s}\le\nu\le\frac{\mu}{c^2},\qquad
 \mu\le\frac{s}{c^2}.
\end{equation}
\end{lemma}

\begin{proof}
The integral expression
\[
 J=\frac12\av{v^2\int_0^1\frac{dt}{1+tv}}
\]
and \(\E J=\lambda^{-2}\E J_e\) imply
\begin{equation}\label{pos:vanishing}
 \E\int_0^1t\av{\frac{v^3}{(1+tv)(1+t\lambda v)}}\,dt=0.
\end{equation}
Indeed, the difference of the two integral expressions for \(J\) and
\(\lambda^{-2}J_e\) is \(-(1-\lambda)/2\) times this integrand.
Define
\[
 A=\E\int_0^1t^2\av{\frac{v^4}{1+tv}}\,dt,
 \qquad
 Q=\frac12\E\int_0^1t^2
       \av{\frac{v^4}{(1+tv)(1+t\lambda v)}}\,dt.
\]
Subtracting \eqref{pos:vanishing} from the corresponding integral
expressions gives the exact formulas
\begin{equation}\label{pos:Q-identities}
 \E J=\mu-\lambda Q,\qquad \eta=A+2\lambda Q.
\end{equation}
For the first, use
\(\mu-\lambda^{-2}\E J_e
 =\frac\lambda2\E\int_0^1 t\av{v^3/(1+t\lambda v)}\,dt\).
For the second, use
\(\eta=\E\int_0^1t\av{v^3}\,dt\).
Both subtractions reduce to multiplication by the two positive
denominators. All these integrals are absolutely integrable: the powers
of \(v\) are bounded, \(1+t\lambda v\ge1-\lambda\), and integration of
the only possibly small denominator \(1+tv\) is controlled by
\(1+|\log m_i|\). Thus the previous logarithmic bounds justify the
interchanges of expectations and integrals.

For \(x>0\), put \(r=\sqrt{x}<1\) and \(k=y/r^3\in[-1,1]\).
Using Equations~\eqref{fnd:cubic} and~\eqref{pos:coordinate-identities},
\begin{equation}\label{pos:h-bound}
 \begin{split}
 h(t)&=x^{-2}\av{\frac{v^4}{1+tv}}
   =\frac{6-10tr k+4t^2r^2k^2}{1-3t^2r^2+2t^3r^3k}\\
 &\ge\frac{6-4tr}{1+tr-2t^2r^2}
  \ge4+\frac{64}{9}(tr-\tfrac12)^2.
 \end{split}
\end{equation}
To verify the first inequality, write \(z=tr<1\). The numerator is
positive and decreases as \(k\) increases to \(1\), while the denominator
increases and remains positive. Evaluation at \(k=1\) gives the displayed
fraction. For the second inequality use
\[
 \frac{6-4z}{1+z-2z^2}
   =4+\frac{8(z-1/2)^2}{1+z-2z^2},\qquad
 0<1+z-2z^2\le\frac98.
\]
Since
\[
 \int_0^1t^2(tr-\tfrac12)^2\,dt
   =\frac{(r-5/8)^2}{5}+\frac1{192},
\]
integration gives \(\int_0^1t^2h(t)\,dt\ge37/27=c_0\).
Consequently \(A\ge c_0\nu\).

At each \(t\), assign the three coordinates weights proportional to
\(v_i^4/(1+tv_i)\). Their weighted mean of \(tv_i\) is
\[
 \frac{\av{v^4}-\av{v^4/(1+tv)}}{\av{v^4/(1+tv)}}
 =\frac6{h(t)}-1\le\frac12.
\]
The function \(z\mapsto(1+\lambda z)^{-1}\) is convex and decreasing
on the relevant interval. Jensen's inequality therefore yields
\[
 \av{\frac{v^4}{(1+tv)(1+t\lambda v)}}
 \ge\frac{2}{2+\lambda}\av{\frac{v^4}{1+tv}}.
\]
Messages with \(x=0\) contribute zero throughout. Integrating proves
\(Q\ge A/(2+\lambda)\ge q\nu\). Now \(J\ge0\) pointwise, because
\(v_i\log(1+v_i)\ge0\). Equations~\eqref{pos:Q-identities}, followed by
Cauchy--Schwarz, prove
\[
 0\le\E J\le\mu-\lambda q\nu,
 \qquad c\nu\le\eta\le\E x^{3/2}\le\sqrt{\mu\nu}.
\]
Finally, Lemma~\ref{fnd:survival} gives \(\mu^2\le s\nu\). Squaring
\(c\nu\le\sqrt{\mu\nu}\) gives \(\nu\le\mu/c^2\), and combining
the two bounds gives the remaining assertion.
\end{proof}

\subsection{Correlation generated by successive branches}

\begin{lemma}\label{pos:correlation-bound}
In either offspring model,
\begin{equation}\label{pos:H-bound}
 \frac{3H}{\alpha}\ge
 \frac{3\mu^2}{2(1+2\mu/3)}-\frac\lambda{\sqrt2}\mu\nu.
\end{equation}
\end{lemma}

\begin{proof}
First consider two independent marginal experiments with posterior means
\(\zeta_0\) and \(\zeta_e\), the latter an edge experiment. Their product
normalizer is \(1+z\), where
\[
 z=\zeta_0\overline{\zeta_e}+\overline{\zeta_0}\zeta_e.
\]
Writing \(\mu_0=\E|\zeta_0|^2\) and
\(\eta_0=\E\Re(\zeta_0^3)\), rotation and reflection symmetry give
\[
 \E z=0,\qquad \E z^2=2\mu_0\lambda^2\mu,
 \qquad \E z^3=2\eta_0\lambda^3\eta.
\]
For \(z>-1\), Taylor's theorem, with fourth derivative
\(2/(1+z)^3>0\), gives
\[
 (1+z)\log(1+z)\ge z+\frac{z^2}{2}-\frac{z^3}{6}.
\]
The inequality extends continuously to \(z=-1\). Thus the correlation
increment \(C_0\), itself a nonnegative relative entropy, satisfies
\begin{equation}\label{pos:increment}
 3C_0\ge3\lambda^2\mu\mu_0-\lambda^3\eta\eta_0.
\end{equation}
The increments telescope to total correlation by the identity for \(I\).

For the regular model, the partial experiment with a fixed number
\(j\le d\) of branches is degraded from the full fixed-point experiment
by discarding branches. For the Poisson model, the partial experiment
is the whole mixture with observed count \(N_t\sim\operatorname{Pois}(t)\),
\(0\le t\le d\); thinning the full experiment produces this mixture.
The degradation comparison in the Poisson case is applied to this
observed-count mixture. Conditional Jensen for the convex function
\(|\zeta|^4\) consequently gives
\(\E|\zeta_0|^4\le\nu\), and hence
\[
 \eta_0\le\sqrt{\mu_0\nu}.
\]
There is also a useful lower bound on \(\mu_0\). Put
\(\beta=\lambda^2\mu\). Each edge posterior mean \(\zeta_e\) has
\[
 \E[\zeta_e\mid S=i]=i\beta,\qquad
 \E[\zeta_e\overline S]=\E|\zeta_e|^2=\beta.
\]
The first identity follows by tilting by the edge likelihood and using
rotation symmetry. For the sum \(T_j\) of \(j\) edge posterior means,
conditional independence therefore gives
\[
 \E[T_j\overline S]=j\beta,\qquad
 \E|T_j|^2=j\beta+j(j-1)\beta^2.
\]
Conditioning on the entire partial observation and applying
Cauchy--Schwarz gives
\begin{equation}\label{pos:projection}
 \mu_j\ge L_j:=\frac{j\beta}{1+(j-1)\beta}\qquad(j\ge1).
\end{equation}
Here \(\mu_j\) denotes its posterior second moment.

We explain carefully how to substitute this lower bound into
\eqref{pos:increment}. Since \(0\le\eta\le\sqrt{\mu\nu}\), any
partial product with \(\mu_0\ge L>0\) satisfies
\[
 3C_0\ge3\lambda^2\mu\mu_0
  \left(1-\frac{\lambda\nu}{3\sqrt{\mu\mu_0}}\right).
\]
The factor in parentheses is increasing in \(\mu_0\). If it is
nonnegative at \(L\), both factors can be bounded below at \(L\).
If it is negative there, the same resulting lower bound follows from
\(C_0\ge0\). In either case,
\begin{equation}\label{pos:truncated-bound}
 3C_0\ge3\lambda^2\mu L-\lambda^3\sqrt\mu\,\nu\sqrt L.
\end{equation}

In the regular case use \(L=L_j\), and note
\(\sqrt{L_j}\le\lambda\sqrt{j\mu}\). Summing the increments yields
\[
 3H\ge3\lambda^4\mu^2\sum_{j=1}^{d-1}
       \frac{j}{1+(j-1)\lambda^2\mu}
       -\lambda^4\mu\nu\sum_{j=1}^{d-1}\sqrt j.
\]
Weighted Jensen, with weights proportional to \(j\), uses
\[
 \lambda^2\frac{\sum_{j=1}^{d-1}j(j-1)}{\sum_{j=1}^{d-1}j}
 =\frac{2(d-2)}{3d}\le\frac23,
 \qquad \lambda^4\sum_{j=1}^{d-1}j=\frac\alpha2.
\]
Cauchy--Schwarz gives
\[
 \frac{\lambda^4}{\alpha}\sum_{j=1}^{d-1}\sqrt j
 \le\frac{\lambda^4}{\alpha}
       \sqrt{(d-1)\frac{d(d-1)}2}
 =\frac\lambda{\sqrt2}.
\]
These are precisely \eqref{pos:H-bound}.

For the Poisson model, let \(N_t\sim\mathrm{Pois}(t)\), \(0\le t\le d\),
and let \(h(t)=\E H_{N_t}\). Because
\(H_n=n\E I_e-\E I_n\) and \(0\le I_n,I_e\le\log3\) in expectation,
\(0\le H_n\le n\log3\). The Poisson series can therefore be
differentiated termwise on bounded intervals, giving
\[
 h'(t)=\E(H_{N_t+1}-H_{N_t}).
\]
Thus \(h'(t)\) is the expected correlation increment on adding one edge
to the partial experiment of intensity \(t\). Indeed, if \(Z_n\) is
the normalizer of the first \(n\) branches and \(Z_{n,e}\) is the
normalizer for their posterior experiment paired with the new edge,
then \(Z_{n+1}=Z_n Z_{n,e}\). Integrating the logarithm of this
identity under the full likelihood tilt gives
\(H_{n+1}-H_n=\E[Z_{n,e}\log Z_{n,e}]\), where the expectation on
the right is under the independent partial and edge marginals.
Averaging over the observed Poisson count gives the asserted increment.
The same projection proof,
now using \(\E N_t=t\) and \(\E N_t(N_t-1)=t^2\), gives
\[
 \mu_t\ge L_t:=\frac{t\beta}{1+t\beta}.
\]
Thinning proves the required degradation and fourth-moment bound. Applying
\eqref{pos:truncated-bound} for \(t>0\), integrating from \(0\) to \(d\),
and using \(h(0)=0\), gives
\[
 3H\ge3\lambda^4\mu^2\int_0^d\frac{t\,dt}{1+t\lambda^2\mu}
        -\lambda^4\mu\nu\int_0^d\sqrt t\,dt.
\]
Weighted Jensen with weights \(t\,dt\) uses
\(\lambda^2\int_0^dt^2\,dt/\int_0^dt\,dt=2/3\), while
\(\lambda^4\int_0^dt\,dt=1/2\). The loss equals
\((2/3)\lambda\mu\nu\), which is at most
\(\lambda\mu\nu/\sqrt2\). As \(\alpha=1\), this proves the claim.
\end{proof}

\subsection{The strict contradiction}

\begin{proposition}\label{pos:nonreconstruction}
At every positive equality parameter \(d\lambda^2=1\), the limiting
posterior is uniform in both tree models. Consequently reconstruction is
impossible for \(0\le\lambda\le d^{-1/2}\).
\end{proposition}

\begin{proof}
Suppose \(\mu>0\). Since \(0\le\E J\le\mu-\lambda q\nu\) and
\(\nu\le\mu/c^2\),
\[
 (\E J)^2\le\mu^2-\lambda q\mu\nu
                  \left(2-\frac{\lambda q}{c^2}\right).
\]
Substituting this estimate and Equation~\eqref{pos:H-bound} into
Equation~\eqref{pos:comparison} gives
\begin{equation}\label{pos:before-scalar}
 \left(\frac{1-\lambda^2}{2\alpha}+g\lambda\mu\right)\nu
  +\frac{3\mu^2}{2(1+2\mu/3)}\le\frac85\mu^2,
 \qquad
 g=\frac85q\left(2-\frac{\lambda q}{c^2}\right)-\frac1{\sqrt2}.
\end{equation}
Here
\[
 q\ge\frac{c_0}{3},\qquad
 \frac{\lambda q}{(c_0+2\lambda q)^2}\le\frac1{8c_0},
\]
and the second inequality is equivalent to
\((c_0-2\lambda q)^2\ge0\). Consequently
\[
 g\ge\frac{113}{81}-\frac1{\sqrt2}>\frac{17}{25}.
\]
For the strict comparison, both sides of
\(113/81-17/25>1/\sqrt2\) are positive, and their squares compare as
\(1448^2/2025^2>1/2\).
Since \(\alpha\le1\) and \(\nu\ge\mu^2/s\),
\eqref{pos:before-scalar} implies
\begin{equation}\label{pos:scalar-necessary}
 T:=\frac{1-\lambda^2}{2s}
       +\frac{3}{2(1+2\mu/3)}
       +\frac{17\lambda\mu}{25s}\le\frac85.
\end{equation}
We show that this is impossible using \(0<s\le1\),
\(\mu\le s/c^2\), and
\begin{equation}\label{pos:c-lower}
 c\ge\underline c(\lambda):=c_0(1+2\lambda/3).
\end{equation}

First suppose \(0<\lambda\le3/10\). Dropping the last term of \(T\)
and using \(\mu\le1/c^2\), it suffices that
\begin{equation}\label{pos:small-sufficient}
 \left(\frac25-\frac{\lambda^2}{2}\right)c^2
 >\frac23\left(\frac{11}{10}+\frac{\lambda^2}{2}\right).
\end{equation}
Indeed, this is exactly the inequality obtained by requiring
\((1-\lambda^2)/2+(3/2)/(1+2/(3c^2))>8/5\).
By \eqref{pos:c-lower}, the left side of \eqref{pos:small-sufficient}
is at least
\[
 c_0^2\left[\frac25+
  \left(\frac8{15}-\frac{3}{20}\left(\frac65\right)^2\right)
       \lambda\right]
 =\frac{2738}{3645}+\frac{162911}{273375}\lambda.
\]
Here we used \(\lambda^2\le(3/10)\lambda\) and
\(1+2\lambda/3\le6/5\), and discarded a nonnegative quadratic term.
The right side is at most \(11/15+\lambda/10\). The constant terms
have difference \(13/729>0\), and
\(162911/273375>1/10\), proving \eqref{pos:small-sufficient}.

Next suppose \(3/10\le\lambda\le1/\sqrt2\). Since
\[
 \frac{3}{2(1+2\mu/3)}\ge\frac32-\mu,
\]
we have
\[
 T-\frac85\ge\frac25-\frac{\lambda^2}{2}
             -\left(1-\frac{17\lambda}{25}\right)\mu
 \ge\frac25-\frac{\lambda^2}{2}
             -\frac{1-17\lambda/25}{c^2}.
\]
Thus it suffices to prove
\[
 f(\lambda):=\left(\frac25-\frac{\lambda^2}{2}\right)
                \underline c(\lambda)^2
 >1-\frac{17\lambda}{25}.
\]
The difference of these two functions is concave, because
\[
 f''(\lambda)=-\frac{1369}{32805}
                   (120\lambda^2+180\lambda+29)<0.
\]
At the first endpoint,
\(f(3/10)=194398/202500>19/20\), whereas
\(1-(17/25)(3/10)=199/250<4/5\). At the other endpoint,
\[
 f(1/\sqrt2)=\frac{1369}{729}
                 \left(\frac{11}{60}+\frac{\sqrt2}{10}\right)
 >\frac35,
 \qquad
 1-\frac{17}{25\sqrt2}<\frac{13}{25}.
\]
The first comparison follows already from \(\sqrt2>7/5\); the second
from \(1/\sqrt2>707/1000\). Concavity now proves the strict inequality
on the entire interval.

Finally suppose \(1/\sqrt2<\lambda<1\), which can occur at equality
only in the Poisson model. Lemma~\ref{fnd:survival}, with
\(d=\lambda^{-2}\), gives \(s\le2(1-\lambda^2)\). Hence
\[
 T-\frac85\ge\frac3{20}
                  -\frac{1-17\lambda/25}{c^2}.
\]
It remains to check
\((3/20)\underline c(\lambda)^2>1-17\lambda/25\).
At \(\lambda=1/\sqrt2\) this is precisely the strict endpoint
inequality just proved. Its left side increases with \(\lambda\), while
its right side decreases, so it holds throughout the remaining interval.

All cases contradict \eqref{pos:scalar-necessary}. Thus \(\mu=0\),
which makes the limiting posterior uniform. Same-sign degradation
extends the conclusion to smaller positive parameters, and
\(\lambda=0\) is the independent channel.
\end{proof}
